% SEGMENT OLP-0663-S01
% Part: proof-theory
% Chapter: natural-deduction
% Section: introduction

\documentclass[../../../include/open-logic-section]{subfiles}

\begin{document}

\olfileid{pt}{nat}{int}

\olsection{அறிமுகம்}

இயல்பான வருவித்தல் என்பது !!{proof} முறைமைகளின் ஒரு குடும்பம்.
இவற்றில் ஒவ்வொரு தருக்க இணைப்பிக்கும் அளவையடைக்கும் இரண்டு அனுமான விதிகள்
உள்ளன: ஒன்று அந்தச் செயலியை ``அறிமுகப்படுத்தும்'' (!!a{introduction} விதி),
மற்றொன்று அதை ``நீக்கும்'' (!!a{elimination} விதி). எடுத்துக்காட்டாக,
\Intro{\lif} விதியின் முடிவு $!A \lif !B$ வடிவிலான !!{formula}s ஆகவும்,
\Elim{\lif} விதியின் முன்கோள்கள் $!A \lif !B$ வடிவிலான
!!{formula}s ஆகவும் இருக்கும்.

% SEGMENT OLP-0663-S02
இயல்பான வருவித்தலின் முதல் இரண்டு வடிவங்களை 1930களில்
கெர்ஹார்ட் ஜென்ட்சனும் (Gerhard Gentzen) ஸ்தானிஸ்வாவ்
யாஷ்கோவ்ஸ்கியும் (Stanisław Jaśkowski) அறிமுகப்படுத்தினர்.
நிறுவல் கோட்பாட்டாளர்கள் பெரும்பாலும் ஜென்ட்சனின் முறைமையையே
விவாதிக்கின்றனர்; கற்பித்தலில் பரவலாகப் பயன்படுத்தப்படும்
முறைமைகளுக்கு யாஷ்கோவ்ஸ்கியின் வடிவம் வழிவகுத்தது. இரண்டு
முறைமைகளிலும் \emph{எடுகோள்கள்}---தாமே !!{prove} செய்யப்பட
வேண்டியதில்லை, ஆனால் பிற !!{formula}s ஐ !!{proving} செய்வதில்
பயன்படும் !!{formula}s---உடன் தொடங்கலாம். முழு !!{proof} உம்
அத்தகைய எடுகோள்களைச் சார்ந்திருக்கலாம். அவை நிறுவப்படாதவை என்பதால்,
அந்த !!{proof} அதன் முடிவை (அதன் \emph{இறுதி-!!{formula}} ஐ)
தனியாக நிறுவாது; எடுகோள்களிலிருந்து முடிவு பின்தொடர்கிறது என்பதையே
நிறுவும்.

% SEGMENT OLP-0663-S03
சில அனுமான விதிகளின் முடிவுகள், முன்னிருந்த சில எடுகோள்களை
இனிச் சார்ந்திருக்காது: அத்தகைய விதிகள் எடுகோள்களை
``விடுவிக்கின்றன''. எடுத்துக்காட்டாக, \Intro{\lif} விதி,
எடுகோள்~$!A$ இலிருந்து முடிவு~$!B$ க்கான !!a{proof} ஐக் கொண்டு
$!A \lif !B$ ஐ முடிவாகத் தருகிறது. $!A \lif !B$ இல் முடியும்
!!{proof} இனி $!A$ ஐச் சார்ந்திருக்காது; எடுகோள் $!A$
விடுவிக்கப்படுகிறது. ஜென்ட்சனின் முறைமைகளில் !!{proof}s என்பது
!!{formula}s ஆல் ஆன மரங்கள்; எடுகோள்கள் அவற்றின் மிக மேலுள்ள
!!{formula}s. \Intro{\lif} போன்ற விதிகளுடன் எந்த எடுகோள்கள்
விடுவிக்கப்படுகின்றன என்பதைக் காட்டும் குறியீடுகள் சேர்க்கப்படுகின்றன.
யாஷ்கோவ்ஸ்கியின் முறைமைகளில் ஓர் எடுகோளைச் சார்ந்த !!{proof} இன்
பகுதி, எடுத்துக்காட்டாக ஒரு செங்குத்துக் கோட்டுக்குப் பின்னால்
உள்தள்ளியோ பெட்டிக்குள்ளோ, தனியே காட்டப்படுகிறது: பெட்டியின்
முதல் வரியில் எடுகோள்~$!A$, இறுதி வரியில் நிபந்தனையின்
பின்விளைவு~$!B$ இருக்கும்; அந்த எடுகோள் விடுவிக்கப்படுகிறது;
\Intro{\lif} இன் முடிவு~$!A \lif !B$ பெட்டிக்கு வெளியே இருக்கும்.

% SEGMENT OLP-0663-S04
இந்த முறைமைகள் ``இயல்பானவை'' எனப்படுவதற்குக் காரணம், அவற்றின்
அனுமான விதிகள் நாம், அல்லது குறைந்தது கணிதவியலாளர்கள்,
இயல்பாகப் பகுத்தறியும் முறையை ஒத்திருப்பதே. ஒரு நிபந்தனைக் கூற்றை
நிறுவ, அதன் முன்னிலைப்பகுதியை எடுகோளாக ஏற்று அதிலிருந்து
பின்விளைவை நிறுவுகிறோம். இதுவே~\Intro{\lif}. நிபந்தனை
$!A \lif !B$ உம் அதன் முன்னிலை~$!A$ உம் தெரிந்தால், $!B$
என்று முடிவுசெய்கிறோம். இதுவே~\Elim{\lif}. ஒரு பிரிப்பிணைவு
$!A \lor !B$ ஐ எடுகோளாகக் கொண்டால் ஒரு முடிவு பின்தொடரும்
என்பதைக் காட்ட, நேர்வுகள் வாரியான நிறுவலைப் பயன்படுத்துகிறோம்;
அதுவே~\Elim{\lor}. பொதுக்கூற்று~$\lforall[x][!A(x)]$ ஐ
நிறுவ, தன்னிச்சையான~$c$ ஒன்றுக்காக $!A(c)$ உண்மை என்பதை
நிறுவுகிறோம்; அதுவே~\Intro{\lforall}. பிற விதிகளும் இவ்வாறே.

% SEGMENT OLP-0663-S05
எடுத்துக்காட்டாக, ஜென்ட்சன் பாணியிலான இயல்பான வருவித்தலில்
$!A \lif (!A \lor !B)$ இன் எளிய !!{proof} இதோ:
\[
\AxiomC{$\Discharge{!A}{x}$}
\RightLabel{\Intro{\lor}}
\UnaryInfC{$!A \lor !B$}
\DischargeRule{\Intro{\lif}}{x}
\UnaryInfC{$!A \lif (!A \lor !B)$}
\DisplayProof
\]
இது எடுகோள்~$!A$ உடன் தொடங்குகிறது; \Intro{\lor} மூலம்
$!A \lor !B$ ஐயும், பின்னர் \Intro{\lif} மூலம்
$!A \lif (!A \lor !B)$ ஐயும் வருவிக்கிறது. \Intro{\lif}
அனுமானம் எடுகோள்~$!A$ ஐ விடுவிக்கிறது; இதைச் சிட்டை~$x$
குறிக்கிறது.

% SEGMENT OLP-0663-S06
நிறுவல் கோட்பாட்டாளர்கள் !!{formula}s ஆல் ஆன மரங்களைக்
கையாளும் ஜென்ட்சனின் மூல முறைமைகளை விரும்புகின்றனர்.
இருப்பினும், அவற்றின் வேறொரு வடிவமும் நிறுவல் கோட்பாட்டில்
முக்கியமானது. எடுகோள்கள்~$\Gamma$ இலிருந்து $!A$ வுக்கான
!!^a{proof}, $!A$ ஆனது $\Gamma$ இலிருந்து பின்தொடர்கிறது,
அதாவது $\Gamma \Entails !A$, என்பதைக் காட்டுகிறது. இயல்பான
வருவித்தலின் அனுமான விதிகளை அத்தகைய பின்தொடர்தல் தொடர்புகளைப்
பற்றிய விதிகளாகவும் வாசிக்கலாம். எடுத்துக்காட்டாக,
$\Gamma + !A \Entails !B$ எனில் $\Gamma \Entails !A \lif !B$
என்பதே \Intro{\lif} விதி கூறுவது. ஆகவே அனுமான விதியின்
முன்கோள்களிலும் முடிவிலும், எடுகோள்களையும் அவற்றிலிருந்து
பின்தொடரும் முடிவையும் நேரடியாகக் கொண்ட தொடரணிகள்
$\Gamma \Sequent !A$ மீது விதிகள் இயங்குமாறு அமைக்கலாம்.
மேலுள்ள எளிய நிறுவல் அத்தகைய முறைமையில் இவ்வாறு இருக்கும்:
\[
\Axiom$x:!A \fCenter !A$
\RightLabel{\Intro{\lor}}
\UnaryInf$x:!A \fCenter !A \lor !B$
\DischargeRule{\Intro{\lif}}{x}
\UnaryInf$\fCenter !A \lif (!A \lor !B)$
\DisplayProof
\]
இங்கும் விதிகள் அவையே: அவை $\Sequent$ க்கு வலப்புறமுள்ள
!!{formula} மீது இயங்குகின்றன; இடப்புறமுள்ள !!{formula}s,
ஒவ்வொரு படியிலும் முடிவு சார்ந்திருக்கும் எடுகோள்களைப் பட்டியலிடுகின்றன.

% SEGMENT OLP-0663-S07
!!{introduction}/!!{elimination} விதிகளின் எல்லா இணைகளும் மட்டும்
செவ்வியல் தருக்கத்திற்குப் போதுமான முழுமையைத் தருவதில்லை.
$\lfalse$ உடன் தொடர்புடைய மேலும் இரண்டு விதிகளைச் சேர்க்க வேண்டும்.
முதல் விதி ``உள்ளுணர்வுவாத அபத்த விதி''~\FalseInt;
இதை ``முரண்பாட்டிலிருந்து எதையும் வருவித்தல்'' என்றும் கூறுவர்:
$\lfalse$ இலிருந்து எதையும் வருவிக்கலாம் என இது சொல்கிறது.
இரண்டாவது, மறைமுக நிறுவலின் செவ்வியல் கொள்கை~\FalseCl{}
(அல்லது ``செவ்வியல் அபத்த விதி''): $\lnot !A$ இலிருந்து
$\lfalse$ ஐ !!{prove} செய்ய முடிந்தால், $!A$ ஐ
!!{prove} செய்ய முடியும். \FalseCl ஐ விலக்கினால்
உள்ளுணர்வுவாதத் தருக்கத்திற்குப் பொருத்தமான முறைமை கிடைக்கும்.
இரண்டையுமே விலக்கினால், அது ``குறைந்தபட்சத் தருக்கம்'' ஆகும்.

% SEGMENT OLP-0663-S08
செவ்வியல் மற்றும் உள்ளுணர்வுவாதத் தருக்கங்களுக்கான
ஜென்ட்சன் பாணி இயல்பான வருவித்தலை இங்கு விவாதிப்போம்.
அவை முறையே \Log{N1c}, \Log{N1i} முறைமைகள். தனித்த
!!{formula}s~$!A$ ஐக் காட்டாமல், தொடரணிகள்
$\Gamma \Sequent !A$ மீது இயங்கும் அவற்றுக்கொத்த
முறைமைகள் \Log{N2c}, \Log{N2i} ஆகும்.

\end{document}
